\documentclass{article}
\usepackage[utf8]{inputenc}
\usepackage{amsmath}
\usepackage{amsfonts}
\usepackage{amssymb}
\usepackage{geometry}
\usepackage{tcolorbox}
\usepackage{hyperref}
\geometry{a4paper, margin=1in}

\title{Module 06: Tree-Based Gradient Boosting Dominance - Part 2}
\author{Cybernauts 2026 Datathon Campaign}
\date{May 2026}

\begin{document}

\maketitle

\begin{abstract}
While LightGBM optimizes for raw speed through asymmetric growth, XGBoost handles dataset volatility with exceptional stability. Known as the legendary algorithm that popularized gradient boosting in competitive data science, XGBoost uses a rigorous optimization engine that balances predictive depth with robust mathematical regularization. This note explores the level-wise growth mechanics of XGBoost, its formalized objective function, and structural regularization strategies.
\end{abstract}

\tableofcontents
\newpage

\section{The Architecture of Level-Wise Growth}
To understand why XGBoost is so stable against distribution noise, we must examine how it builds individual decision boundaries compared to leaf-wise frameworks.

\subsection{Level-Wise (Depth-Wise) Splitting}
Unlike LightGBM's leaf-wise strategy, XGBoost typically implements a **Level-Wise (Depth-Wise) Growth** engine. When a tree expands, the algorithm evaluates splits across an entire horizontal layer simultaneously. It completes a uniform layer of depth before moving down to compute subsequent sub-splits.

\textit{The Stability Advantage:} Level-wise tree generation maintains a balanced, symmetric structure. This structural uniformity acts as a built-in constraint. Because it cannot create isolated, deep branches based on small clusters of data points, it naturally resists overfitting localized noise, making it highly reliable across varied data distributions.

\section{The Regularized Objective Function}
The defining feature of XGBoost is its mathematical formulation. It embeds regularization directly into its optimization goal, penalizing model complexity at every step of the learning process.

\subsection{The Formal Math}
At step $t$ of the boosting process, the objective function to minimize is formalized as:
\begin{equation}
    \mathcal{O}^{(t)} = \sum_{i=1}^{n} l\left(y_i, \hat{y}_i^{(t-1)} + f_t(x_i)\right) + \Omega(f_t)
\end{equation}
Where $l$ represents the differentiable loss function between targets $y_i$ and current predictions, and $\Omega(f_t)$ is the explicit complexity penalty for the new tree $f_t$.

The structural penalty $\Omega(f_t)$ targets both the leaf configuration and the output weights:
\begin{equation}
    \Omega(f_t) = \gamma T + \frac{1}{2}\lambda \sum_{j=1}^{T} w_j^2 + \alpha \sum_{j=1}^{T} |w_j|
\end{equation}
Where $T$ is the number of leaves, $w_j$ represents the continuous output score weights of leaf $j$, $\lambda$ is the L2 regularization coefficient, and $\alpha$ is the L1 regularization coefficient.

\section{Core Hyperparameters to Master}
Controlling the balance between gradient alignment and structural complexity depends on three main tuning handles:

\subsection{\texttt{max\_depth}: The Structural Structural Ceiling}
This parameter dictates the absolute maximum depth a decision tree can reach.
\begin{itemize}
    \item **The Operational Rule:** Because XGBoost grows level-by-level, increasing \texttt{max\_depth} expands the tree exponentially, creating $2^{\text{max\_depth}}$ potential leaves. In tabular hackathons, \texttt{max\_depth} is typically constrained between $3$ and $8$ to keep individual trees shallow and generalizable.
\end{itemize}

\subsection{\texttt{reg\_lambda} ($\lambda$): The L2 Weight Smoothing Shield}
The L2 regularization term applied to the leaf weights. It corresponds directly to the square of the $L_2$ vector norm.
\begin{itemize}
    \item **The Operational Rule:** Increasing \texttt{reg\_lambda} forces leaf weights toward zero, smoothing out the model's predictions. It is your primary defense against volatile feature matrices and large outliers, as it prevents any single leaf node from making extreme, overconfident predictions.
\end{itemize}

\subsection{\texttt{reg\_alpha} ($\alpha$): The L1 Feature Sparsity Driver}
The L1 regularization term applied to the leaf weights. It corresponds directly to the absolute $L_1$ vector norm.
\begin{itemize}
    \item **The Operational Rule:** Unlike L2 regularization, increasing \texttt{reg\_alpha} can drive leaf weights exactly to zero, introducing structural sparsity. This helps drop uninformative, noisy feature combinations completely, keeping the overall ensemble stream streamlined and focused on high-signal patterns.
\end{itemize}

\newpage
\section{Practice Problems}

\textbf{P1:} Explain how XGBoost's level-wise tree growth acts as an implicit regularizer compared to LightGBM's leaf-wise tree growth when dealing with a tabular feature filled with extreme sample outliers.
\\
\textit{Answer: LightGBM's leaf-wise engine focuses entirely on maximizing loss reduction, meaning it can create an asymmetric branch that digs deeply into a small cluster of extreme outliers to isolate them. XGBoost's level-wise growth avoids this by forcing splits to occur uniformly across an entire horizontal layer. To isolate a small pocket of outliers, XGBoost would have to split across all other nodes at that same level as well. The internal objective function will block this expansion if the global gain doesn't justify the complexity penalty, keeping the model from overfitting to outlier noise.}

\newline
\textbf{P2:} Review the regularization penalty equation: $\Omega(f) = \gamma T + \frac{1}{2}\lambda \sum w_j^2 + \alpha \sum |w_j|$. Prove mathematically what happens to the optimal leaf weight $w_j^*$ as the L2 regularization hyperparameter $\lambda$ approaches infinity ($\infty$).
\\
\textit{Answer: The optimal leaf weight solution for a given tree node is derived by setting the partial derivative of the objective function with respect to $w_j$ to zero, which yields:
\begin{equation*}
    w_j^* = -\frac{G_j}{H_j + \lambda}
\end{equation*}
Where $G_j$ and $H_j$ represent the sums of the first and second-order loss gradients within leaf $j$. Taking the limit as $\lambda \to \infty$:
\begin{equation*}
    \lim_{\lambda \to \infty} w_j^* = \lim_{\lambda \to \infty} \left( -\frac{G_j}{H_j + \lambda} \right) = 0
\end{equation*}
As $\lambda$ approaches infinity, the output leaf weights are crushed down to exactly $0$, transforming the tree into a flat, neutral predictor that outputs zero change.}

\newline
\textbf{P3:} You are optimizing an XGBoost model. Your training pipeline runs smoothly, but your local 5-Fold validation loop shows severe overfitting: your training log-loss is $0.112$ while your out-of-fold validation log-loss is $0.425$. Suggest a clear hyperparameter tuning strategy using \texttt{max\_depth}, \texttt{reg\_lambda}, and \texttt{reg\_alpha} to fix this gap.
\\
\textit{Answer: To close the gap and prevent overfitting, you need to reduce model complexity and increase regularization penalties:
\begin{enumerate}
    \item **Decrease** \texttt{max\_depth} (e.g., drop it from 8 down to 4 or 5) to limit how deep individual trees can grow.
    \item **Increase** \texttt{reg\_lambda} (e.g., raise it from its default of 1 to 5 or 10) to smooth out the leaf weights and reduce prediction variance.
    \item **Increase** \texttt{reg\_alpha} (e.g., raise it from 0 to 1 or 2) to introduce sparsity and help drop uninformative feature noise.
\end{enumerate}}

\end{document}
